% ============================================================================
% 第 4 章 问题形式化建模与物理约束（待与队友商议合写）
% ============================================================================
\section{问题形式化建模与物理约束}
\label{sec:modeling}

\begin{planbox}{第 4 章写作规划（供双人分工确认）}
  本章拟以严谨的运筹学与多智能体路径规划（MAPF）记号给出全系统的统一数学表述，建议细分为以下小节并在商定分工后补充完整公式与示意配图：
  \begin{enumerate}
    \item \textbf{4.1 仓库拓扑图与多层货架建模}：定义四邻接离散栅格图 $G=(V,E)$、通道/货架/立柱/交付站/停靠区集合，分别给出主线 A（Legacy / V4 单面多层 / V9.1 混合背靠背货架）与主线 B（$62\times 47$、6 车、5 层 WEPAStacks 映射网格）的布局示意图。
    \item \textbf{4.2 订单任务与生命周期状态机}：形式化描述订单发布时刻 $r_j$、SKU 与货架层高 $l_j$、入/出库任务重叠规则、装货持续时间（如 V4 的 $2 l_j$ 步）、交付站卸货即离、波次间安全待命及流关闭（\code{close}）后的最终返库状态转移。
    \item \textbf{4.3 时空无碰撞与仓储工艺硬约束}：列出离散步长下顶点互斥约束、相向边交换互斥约束、连续装卸占位不变性、同面货架层序解锁偏序关系以及在线调度的因果信息边界条件。
    \item \textbf{4.4 评价指标与目标函数定义}：给出完工时间 $T$、平均响应时间 $R$、含返库总里程 $D$、主线 B 固定尺度线性目标 $F_{\text{WEPA}}$、主线 A 参照归一化目标 $F_{\text{Core}}$ 以及统一改善率公式的严格定义；并简要说明微型两车 CP-SAT 精确参照模型的作用。
  \end{enumerate}
\end{planbox}
